% This file is generated from the corresponding Obsidian note.
% Source: Teoricos/GS - Teo7.md
\chapter{Fibrado tangente y diferencial global}
\label{ch:teo-7}
\section{La diferencial de una función en coordenadas}
\begin{bookproposition}{Matriz de la diferencial en coordenadas}
Sean $M$ y $N$ variedades suaves de dimensiones $m$ y $n$, respectivamente, y sea $F:M\to N$ suave. Fijemos $p\in M$, escribamos $q=F(p)$ y tomemos cartas suaves $(U,\varphi=(x_1,\dots,x_m))$ al rededor de $p$ y $(V,\psi=(y_1,\dots,y_n))$ al rededor de $q$ con $F(U)\subseteq V$. Si se consideran las bases coordenadas del $T_{p}M$ y $T_{q}N$ respectivamente

\[
\mathcal B_M=\left\{\left.\frac{\partial}{\partial x_1}\right|_p,\dots,\left.\frac{\partial}{\partial x_m}\right|_p\right\},\qquad\mathcal B_N=\left\{\left.\frac{\partial}{\partial y_1}\right|_q,\dots,\left.\frac{\partial}{\partial y_n}\right|_q\right\}
\]
entonces la matriz de $dF_p:T_pM\to T_qN$ respecto de estas bases es la jacobiana de

\[
\widehat F:=\psi\circ F\circ\varphi^{-1}:\varphi(U)\subseteq\mathbb R^m\to\mathbb R^n
\]
evaluada en $\varphi(p)$:

\[
[dF_p]_{\mathcal B_M}^{\mathcal B_N}=J\widehat F(\varphi(p))
\]
\end{bookproposition}
\begin{bookproof}{}
\begin{enumerate}
\item Escribimos
\end{enumerate}
\[
dF_p\left(\left.\frac{\partial}{\partial x_j}\right|_p\right)=\sum_{i=1}^n c_{ij}\left.\frac{\partial}{\partial y_i}\right|_q
\]
\begin{enumerate}[start=2]
\item Aplicando ambos lados a la función coordenada $y_i$ se obtiene
\end{enumerate}
\[
c_{ij}=dF_p\left(\left.\frac{\partial}{\partial x_j}\right|_p\right)(y_i) = \left.\frac{\partial}{\partial x_j}\right|_p(y_i\circ F)
\]
\begin{enumerate}[start=3]
\item Como $y_i=\pi_i\circ\psi$ resulta
\end{enumerate}
\[
y_i\circ F=\pi_i\circ\psi\circ F=\pi_i\circ\widehat F\circ\varphi
\]
y por definición de vector coordenado

\[
c_{ij}=\left.\frac{\partial}{\partial x_j}\right|_{\varphi(p)}(\pi_i\circ\widehat F)=\frac{\partial \widehat F_i}{\partial x_j}(\varphi(p))
\]
\begin{enumerate}[start=4]
\item Luego la matriz $(c_{ij})$ coincide exactamente con $J\widehat F(\varphi(p))$.
\end{enumerate}
\bookqed
\end{bookproof}
\phantomsection\label{blk:bb7fdf}
\section{Vectores tangentes como velocidades de curvas}
\begin{bookdefinition}{Curva suave}
Una función suave

\[
\alpha:(a,b)\subseteq\mathbb R\to M
\]
se llama \textbf{curva suave} en $M$.

\end{bookdefinition}
\begin{bookdefinition}{vector tangente a una curva}
Sea $\alpha:(a,b)\to M$ una curva suave y sea $t\in(a,b)$. El vector tangente a $\alpha$ en el instante $t$ es

\[
\alpha'(t):=d\alpha_t\left(\left.\frac{\partial}{\partial s}\right|_t\right)\in T_{\alpha(t)}M,
\]
donde $(\mathbb R,\operatorname{id}=s)$ es la carta canónica de $\mathbb R$.

\end{bookdefinition}
\phantomsection\label{blk:0106aa}
\begin{bookexercise}{Curvas y vectores tangentes}
Sea $M$ una variedad diferenciable, $I\subseteq\mathbb R$ un intervalo y $\gamma:I\to M$ una curva diferenciable.

\begin{enumerate}
\item Mostrar que efectivamente $\gamma'(t)\in T_{\gamma(t)}M$.
\item Probar que
\end{enumerate}
\[
\gamma'(t)=(d\gamma)_t\left(\left.\frac{\partial}{\partial s}\right|_t\right).
\]
\begin{enumerate}[start=3]
\item Si $(U,\varphi=(x_1,\dots,x_n))$ es un sistema coordenado alrededor de $\gamma(0)$ y
\end{enumerate}
\[
(\varphi\circ\gamma)(t)=(r_1(t),\dots,r_n(t)),
\]
mostrar que

\[
\gamma'(0)=\sum_{i=1}^n r_i'(0)\left.\frac{\partial}{\partial x_i}\right|_{\gamma(0)}.
\]
\begin{enumerate}[start=4]
\item Mostrar que para todo $p\in M$ y todo $v\in T_pM$ existe una curva suave $\sigma:(-\varepsilon,\varepsilon)\to M$ tal que
\end{enumerate}
\[
\sigma(0)=p,\qquad \sigma'(0)=v.
\]
\end{bookexercise}
\begin{bookcorollary}{La diferencial lleva velocidades en velocidades}
Sea $F:M\to N$ una función suave, $p\in M$ y $v\in T_pM$. Si $\alpha:(-\varepsilon,\varepsilon)\to M$ es una curva suave con

\[
\alpha(0)=p,\qquad \alpha'(0)=v,
\]
y definimos

\[
\beta:=F\circ\alpha
\]
entonces

\[
(dF)_p(v)=\beta'(0)
\]
\end{bookcorollary}
\begin{bookproof}{}
\begin{enumerate}
\item Por la regla de la cadena
\end{enumerate}
\[
d(F\circ\alpha)_0=dF)_p\circ(d\alpha)_0
\]
\begin{enumerate}[start=2]
\item Aplicando ambos lados al vector $\left.\frac{\partial}{\partial s}\right|_0\in T_0\mathbb R$ y usando la definicion \hyperref[blk:0106aa]{Teórico 7} se obtiene
\end{enumerate}
\[
\beta'(0)=d(F\circ\alpha)_0\left(\left.\frac{\partial}{\partial s}\right|_0\right)=(dF)_p\left(d\alpha_0\left(\left.\frac{\partial}{\partial s}\right|_0\right)\right)=(dF)_p(v)
\]
\bookqed
\end{bookproof}
\section{El fibrado tangente}
\begin{bookdefinition}{Fibrado tangente}
Sea $M$ una variedad suave. Se define el \textbf{fibrado tangente} de $M$, denotado por $TM$, como la unión disjunta de todos los espacios tangentes:

\[
TM:=\bigsqcup_{p\in M}T_pM
=
\bigcup_{p\in M}\{(p,v):v\in T_pM\}.
\]
La aplicación

\[
\pi:TM\to M,\qquad (p,v)\mapsto p
\]
se llama \textbf{proyección natural}.

\end{bookdefinition}
\phantomsection\label{blk:4762ec}
\begin{bookremark}{Notación}
Dado un punto $(p,v)\in TM$, suele escribirse simplemente $v$ en lugar de $(p,v)$ cuando el punto base $p$ está claro por el contexto. O si no $v_{p}$

\end{bookremark}
\begin{booktheorem}{$TM$ es una variedad suave}
Si $M$ es una variedad suave de dimensión $n$, entonces $TM$ tiene una topología natural y una estructura suave que lo convierten en una variedad suave de dimensión $2n$. Además, la proyección natural

\[
\pi:TM\to M
\]
es suave.

\end{booktheorem}
\begin{bookproof}{}
\begin{enumerate}
\item Primero demos cartas que definan la estructura sobre $TM$
\item Sea $(U,\varphi=(x_1,\dots,x_n))$ una carta suave de $M$ y sea
\end{enumerate}
\[
\widetilde U:=\pi^{-1}(U)\subseteq TM
\]
\begin{enumerate}[start=3]
\item Automaticamente tenemos base de $T_{p}M$ osea dado $v\in T_pM$ con $p\in U$ se escribe de manera única como
\end{enumerate}
\[
v=\sum_{i=1}^n v_i\left.\frac{\partial}{\partial x_i}\right|_p
\]
donde $v_i=v(x_i)$. 4. Definimos

\[
\widetilde\varphi:\widetilde U\subseteq TM\to\varphi(U)\times\mathbb R^n\subseteq\mathbb R^{2n}
\]
por

\[
\widetilde\varphi(p,v)=\bigl(x_1(p),\dots,x_n(p),v_1,\dots,v_n\bigr)
\]
llamada \textbf{carta inducida} 5. Esta aplicación es biyectiva, y su inversa es

\[
\widetilde\varphi^{-1}(r_1,\dots,r_n,t_1,\dots,t_n)=\left(\varphi^{-1}(r_1,\dots,r_n),\sum_{i=1}^n t_i\left.\frac{\partial}{\partial x_i}\right|_{\varphi^{-1}(r_1,\dots,r_n)}\right).
\]
\begin{enumerate}[start=6]
\item Esta claro que estas cartas cubren todo $TM$. Veamos que son compatibles
\item Sean $(\tilde{V},\tilde{\psi})$ y $(\tilde{U},\tilde{\varphi})$ dos cartas de $TM$ asociadas a $(V,\psi=(y_1,\dots,y_n))$ y $(U,\varphi=(x_{1},\ldots x_{n}))$ tales que $\tilde{V}\cap \tilde{U}\neq 0$, luego $V\cap U\neq0$
\item Entonces
\end{enumerate}
\[
\tilde{\psi}\circ \tilde{\varphi}^{-1}:\tilde{\varphi}(\tilde{U}\cap \tilde{V}) \subset \mathbb{R}^{2n}\longrightarrow\tilde{\psi}(\tilde{U}\cap \tilde{V}) \subset \mathbb{R}^{2n}
\]
que es:

\[
\begin{aligned}(r_{1},\ldots r_{n},t_{1},...,t_{n})&\longmapsto^{\tilde{\varphi}^{-1}} \left(\tilde{\varphi}^{-1}(r),\sum_{i=1}^n t_i\left.\frac{\partial}{\partial x_i}\right|_{\varphi^{-1}(r)}\right)\\&\longmapsto^{\tilde{\psi}}(y_{1}(\tilde{\varphi}^{-1}(r)),\ldots,y_{n}(\tilde{\varphi}^{-1}(r)),\sum_{i=1}^n t_i\left.\frac{\partial}{\partial x_i}\right|_{\varphi^{-1}(r)}y_{1},\ldots,\sum_{i=1}^n t_i\left.\frac{\partial}{\partial x_i}\right|_{\varphi^{-1}(r)}y_{n}) \\&=\bigg((\psi\circ\varphi ^{-1})(r),\sum_{i=1}^n t_i\left.\frac{\partial}{\partial r_i}\right|_{r}(y_{1}\circ\varphi ^{-1}),\ldots,\sum_{i=1}^n t_i\left.\frac{\partial}{\partial r_i}\right|_{r}(y_{n}\circ\varphi ^{-1})\bigg)\\&=\bigg((\psi\circ\varphi ^{-1})(r),\sum_{i=1}^n t_i\frac{\partial}{\partial r_i}(\psi\circ\varphi ^{-1})_{1}(r),\ldots,\sum_{i=1}^n t_i\frac{\partial}{\partial r_i}(\psi\circ\varphi ^{-1})_{n}(r)\bigg)\end{aligned}
\]
\begin{enumerate}[start=9]
\item Obviamente las primeras coordendas son suaves por que son el cambio de coordenadas, y la segunda tanda de coordenadas cada derivada parcial es suave por que $\psi\circ\varphi ^{-1}$ es suave y proyectar en la primera coordenada es suave entonces son sumas de suaves
\item Y obviamente cambiando los roles obtenemos que la otra composicion es suave tambien
\item Por lo tanto, las cartas $\{(\widetilde U,\widetilde\varphi)\}$ formarian un atlas suave sobre $TM$ de dimensión es $2n$. (Si $\tilde{U},\tilde{V}$ fueran abiertos)
\item Sea
\end{enumerate}
\[
\mathcal{B}=\{\tilde{\varphi}^{-1}(W): W \text{ es abierto de } \mathbb{R}^{2n},\ (U,\varphi) \text{ es carta suave de } M\}
\]
\begin{enumerate}[start=13]
\item Veamos que $\mathcal{B}$ cumple las condiciones para ser una base de una topología (para $TM$)
\item $\mathcal{B}$ cubre $TM$. Pues el conjunto de cartas suaves de $M$ cubre a $M$, y en $\mathcal{B}$ están los $\tilde{U}=\pi^{-1}(U)$, con $U$ dominio de carta suave.
\item Sean $\tilde{\varphi}_1^{-1}(W_1)$ y $\tilde{\varphi}_2^{-1}(W_2)$ en $\mathcal{B}$, y sea $z \in \tilde{\varphi}_1^{-1}(W_1)\cap \tilde{\varphi}_2^{-1}(W_2)$. Queremos ver que existe un miembro de $\mathcal{B}$ que contiene a $z$ y está contenido en la intersección.
\item Para ver esto, alcanza con ver que $\tilde{\varphi}_1^{-1}(W_1)\cap \tilde{\varphi}_2^{-1}(W_2)\in \mathcal{B}$. Llamemos $\tilde{U}_1=\operatorname{dm}\tilde{\varphi}_1$ y $\tilde{U}_2=\operatorname{dm}\tilde{\varphi}_2$. Entonces
\end{enumerate}
\[
\tilde{\varphi}_1^{-1}(W_1)\cap \tilde{\varphi}_2^{-1}(W_2)=\tilde{\varphi}_1^{-1}(W_1)\cap \tilde{\varphi}_{1}^{-1}\circ\tilde{\varphi}_1\circ \tilde{\varphi}_{2}^{-1}(W_{2}\cap \tilde{\varphi}_2(\tilde{U}_{1}\cap \tilde{U}_{2}))
\]
\begin{enumerate}[start=17]
\item Como las funciones biyectivas se llevan bien con las uniones e intersecciones
\end{enumerate}
\[
\tilde{\varphi}_1^{-1}(W_1)\cap \tilde{\varphi}_2^{-1}(W_2)=\tilde{\varphi}_1^{-1}(W_1\cap \tilde{\varphi}_{1}\circ \tilde{\varphi}_{2}^{-1}(W_{2}\cap \tilde{\varphi}_2(\tilde{U}_{1}\cap \tilde{U}_{2}))
\]
\begin{enumerate}[start=18]
\item Y $W_{1}\cap \tilde{\varphi}_{1}\circ \tilde{\varphi}_{2}^{-1}(W_{2}\cap \tilde{\varphi}_{2}(\tilde{U_{1}}\cap \tilde{U}_{2}))$ es abierto por que $W_{1}$ es abierto.
\item Y ademas vimos en 8. y 9. que $\tilde{\varphi}_{1}\circ\tilde{\varphi}_{2}^{-1}$ es difeo entre sus abiertos $\tilde{\varphi}_{1}(\tilde{U_{1}}\cap \tilde{U_{2}})$ y $\tilde{\varphi}_{2}(\tilde{U_{1}}\cap \tilde{U_{2}})$ por ende es homeomorfismo luego $W_{2}\cap \tilde{\varphi_{2}}(\tilde{U}_{1}\cap \tilde{U}_{2})$ es abierto de $\mathbb{R}^{2n}$ por ser interseccion de abiertos
\item Mostrando que $\tilde{\varphi}_1^{-1}(W_1)\cap \tilde{\varphi}_2^{-1}(W_2)\in \mathcal{B}$ como queriamos
\item Esta topología tiene una base numerable. Tomamos una base numerable de dominios de cartas de $M$, digamos $\{U_i\}$, y entonces $\{\tilde{U}_i=\pi^{-1}(U_i)\}$ es una base numerable para $TM$.
\item ¿Por qué se puede hacer esto? Porque $M$ es $N2$ por definicion de variedad topologica luego existe una base numerable formada por dominios de cartas, y la proyección $\pi:TM\to M$ permite levantar esos abiertos.
\item Chequeemos que $TM$ es Hausdorff. Sean $(p,v)\neq(q,w)$ en $TM$. Tenemos dos casos:
\end{enumerate}
\begin{itemize}
\item Si $p\neq q$

\begin{enumerate}
\item como $M$ es Hausdorff, existen abiertos disjuntos $U,V\subset M$ tales que $p\in U$ y $q\in V$. Entonces $\tilde{U}=\pi^{-1}(U)$ y $\tilde{V}=\pi^{-1}(V)$ son abiertos disjuntos por \hyperref[blk:4762ec]{Teórico 7} y que separan $(p,v)$ y $(q,w)$.
\end{enumerate}
\item Si $p=q$

\begin{enumerate}
\item Entonces $v\neq w$. Tenemos que $(p,v),(p,w)\in \tilde{U}$ con $(U,\varphi=(x_{1},\ldots x_{n}))$ cartas de $M$ en $p$.
\item Escribimos
\end{enumerate}
\end{itemize}
\[
v=\sum t_i\frac{\partial}{\partial x_{i}}\big|_p \qquad w=\sum s_i\frac{\partial}{\partial x_{i}}\big|_p
\]
\begin{Verbatim}[fontsize=\small]
 3. Como $v\neq w$, se tiene $(t_1,\dots,t_n)\neq(s_1,\dots,s_n)$.
 4. Luego existen abiertos disjuntos $B_1,B_2\subset \mathbb{R}^n$ que separan estos puntos. 
 5. Identificando $\mathbb{R}^{2n}\cong\mathbb{R}^n\times\mathbb{R}^n$, Tenemos que 
\end{Verbatim}
\[
W_{1}=\varphi(U)\times B_1 \qquad \varphi(U)\times B_2=W_2
\]
son abiertos de $\mathbb{R}^{2n}$ que contienen a $\tilde{\varphi}(p,v)$ y $\tilde{\varphi}(p,w)$ y los separan. 6. Por lo tanto, $\tilde{\varphi}^{-1}(W_1)$ y $\tilde{\varphi}^{-1}(W_2)$ son abiertos de $\mathcal{B}$ que separan a $(p,v)$ y $(p,w)$. 7. Por último, la proyección $\pi:TM\to M$ es suave. Esto se verifica usando las cartas inducidas $\{(\tilde{U},\tilde{\varphi})\}$ asociadas a $\{(U,\varphi)\}$ de $M$ . 8. En coordenadas se tiene $\varphi\circ\pi\circ\tilde{\varphi}^{-1}:\tilde{\varphi}(\tilde{U})\subset\mathbb{R}^{2n}\to\mathbb{R}^n$, $(x,v)\mapsto x$, que es claramente una aplicación suave.

\bookqed
\end{bookproof}
\begin{bookremark}{Fórmula de los cambios de coordenadas en $TM$}
Si

\[
\widetilde\psi\circ\widetilde\varphi^{-1}(r,t)=\bigl(r',t'\bigr),
\]
entonces

\[
r'=\psi\circ\varphi^{-1}(r),
\qquad
t'=J(\psi\circ\varphi^{-1})(r)\,t.
\]
Es decir, las coordenadas de base cambian por la función de transición usual, y las coordenadas de la fibra cambian por su jacobiana.

\end{bookremark}
\begin{bookexercise}{Caso de una sola carta}
Si $M$ puede cubrirse con una sola carta suave, entonces

\[
TM\cong M\times\mathbb R^n
\]
como variedades suaves.

\end{bookexercise}
\begin{bookexample}{Opcional}
Sea $M$ una variedad que puede ser cubierta con una sola carta, entonces $TM$ es difeomorfa a $M\times \mathbb{R}^{n}$

\end{bookexample}
\section{Diferencial global}
\begin{bookdefinition}{Diferencial global}
Sea $F:M\to N$ una función suave. Se define la \textbf{diferencial global} de $F$ como la aplicación

\[
dF:TM\to TN,
\]
dada por

\[
dF(p,v)=\bigl(F(p),(dF)_p(v)\bigr).
\]
Esta aplicación extiende lo que hace $(dF)_{q}$ con $q\in M$ a toda la variedad.

\end{bookdefinition}
\phantomsection\label{blk:fce8ee}
\begin{bookremark}{Suavidad de la diferencial global}
Mostrar que si $F:M\to N$ es suave, entonces

\[
dF:TM\to TN
\]
también es suave.

\end{bookremark}
\begin{bookproof}{}
Se ve en el practico

\bookqed
\end{bookproof}
\phantomsection\label{blk:a36c7a}
\begin{bookremark}{El funtor tangente Lee pagina 73}
Si $\mathbf{Diff}$ denota la categoría cuyas:

\begin{itemize}
\item objetos son las variedades suaves,
\item morfismos son las funciones suaves,
\end{itemize}
entonces la asignación

\[
(M,F)\longmapsto (TM,dF)
\]
define un \textbf{funtor covariante}

\[
T:\mathbf{Diff}\to\mathbf{Diff},
\]
llamado \textbf{funtor tangente}.

\end{bookremark}
